% !TeX root = ../Masterarbeit.tex

\documentclass[Masterarbeit.tex]{subfile}

\begin{document}
	
	\chapter{Data and Methods}
	
	\section{Data}
	
	In order to gain a new estimate of North Atlantic ocean heat content, this work will use two different sources of training and testing data for the utilized U-net. On the one hand, data assimilation runs from the Max Planck Institute Earth System Model (MPI-ESM, low resolution, \cite{mauritsen_developments_2019}) are deployed to train the neural network to reconstruct temperature maps. On the other hand, observational data from the fourth version of the Met Office Hadley Center 'EN' series (EN4, \cite{good_en4_2013}) subsurface ocean temperature profiles are applied as basis for the final reconstruction. After the neural network reconstructions, the temperature profiles are reverted back to total values and converted into a two dimensional OHC map via integration:
	
	
	\begin{align}
	H = c_{p} \int_{z_0}^{z_1} \rho_z \cd T_z dz ,
	\end{align}
	
	In this case $H$ denotes the ocean heat content, $c_{p}$ the specific heat capacity of sea water, $\rho_z$ the density of the water at depth $z$ and $T_z$ the temperature at this depth. This OHC field is then reduced to one value by summing over all relevant gridpoints for the calculation of an OHC timeseries.	
	
	
	\subsection{Data Assimilation}
	
	In general, Kalman Filter (KF) assimilations assume the uncertainties of model and observations to be Gaussian and then assimilates the observations sequentially. Both, model and observations, are described as probability density functions (PDFs) and the Kalman filter alternates between forecasting and analyzing the PDF's covariance and mean \cite{carrassi_data_2018}. It is often supplemented by a backward calculation of the smoothing mean and covariance to refine the previous estimates in what is called the Kalman smoother (KS). By introducing an ensemble of model trajectories instead of one single representation of the model dynamics, the KF can be adjusted and extended in what is called an ensemble Kalman filter (EnKF). This allows for the covariance matrices to be approximated by an ensemble of model states. The covariance matrix and ensemble mean are then sequentially forecast and updated. EnKFs provide computationally efficient and very accurate results, especially in high dimensional systems, such as the climate and Earth system \citep{carrassi_data_2018}. \\
		
	For this work, an ensemble coupled MPI-ESM assimilation \citep{brune_mpi-esm-lr_1201p5_2021} is used to train the U-net to reconstruct masked temperature profiles. An oceanic localized EnKF, implemented with the Parallel Data Assimilation Framework \citep{nerger_software_2013}, assimilates monthly EN4 temperature and salinity profiles. Simultaneously, atmospheric nudging of ERA40/ERAInterim/ERA5 reanalyses \citep{uppala_era-40_2005, dee_era-interim_2011, hersbach_era5_2020} takes place. All monthly datasets range from January 1958 until October 2020 and present temperature profiles up to a depth of 2.000 m in 40 depth steps.\\
	These assimilation runs are based on the work of \cite{brune_assimilation_2015} and use a version of the singular evolutive interpolated Kalman filter (SEIK, \cite{pham_filtres_1998}) in their EnKF. This assimilation was developed as a model inherent initialization for the MPI-ESM and showed improvements especially in the assimilation of sea surface temperatures and ocean heat content \citep{brune_assimilation_2015}. Therefore, it provides a good estimate of the relevant variables and represents a fitting choice as training data for the U-net.
	
	
	\subsection{Observational Data}
	
	The EN4 project describes a collection of global ocean temperature and salinity profile datasets spanning from 1900 until present day. It contains observational data from Argo floats, Arctic Synoptic Basinide Oceanography (ASBO), the Global Temperature and Salinity Profile Program (GTSPP) and the World Ocean Database (WOD13) \citep{good_en4_2013}. As these datasets partly overlap, input data is preprocessed and duplicates sorted out before the complete EN4 dataset is created. The same procedure is applied to profiles that are not necessarily identified to be the same but that are very close to each other in time and space \citep{good_en4_2013}. Additionally, various quality control checks are performed in order to identify data profiles with faulty location or measurement data, such as missing depth steps or the location assignment to a point on land rather than in the ocean. All remaining datapoints are then collected in NetCDF files and are publicly available under \url{https://www.metoffice.gov.uk/hadobs/} \citep{good_en4_2013}. Furthermore, the EN4 project also provides an interpolated version of the observed data points in the form of a reanalysis. These are created by combining a persistence based forecast from the previous month's reanalysis with the observed temperature profiles \cite{good_en4_2013}. As the persistence relaxes towards the baseline climatology when not corrected by observations, the EN4 reanalysis tends to posses a lot of uncertainty in times with low observational density. The applied baseline climatology is taken from the years 1970 -- 2000. Thus, the overall cooling trend in the NA during that time (\ref{fig:ensemble_spread}) leads to a cool bias especially in the EN4 OHC estimate especially in the remaining pre-Argo era \cite{good_en4_2013}. Despite these shortcomings, the EN4 will be taken as a reference point in the evaluation of the network reconstructions as they are highly consistent with the observational data, especially in the well observed Argo era.
	
	\section{Methods}
	
	\subsection{The Convolutional U-net} \label{method:unet}
	
	The U-net adopted in this work is based on the one used by \cite{kadow_artificial_2020} to reconstruct near-surface air temperature data. However, a few modifications have been made for it to fit the task of reconstructing ocean heat content data instead. In general, the U-net is set up to incorporate four encoding and four max pooling layers. These are accompanied by the same amount of decoding and upsampling layers on the decoding path as depicted in figure \ref{fig:unet}. Each encoding and decoding layer consists of a partial convolution of the input image. Partial convolutions use a masked and re-normalized convolution operation, which is then followed by a mask-update step \citep{liu_image_2018}. In mathematical terms, this means that the partial convolution at each location is defined as
	
	\begin{align}
		x' = \begin{cases}
		W^{T} (X \cdot M) \frac{\sum(1)}{\sum(M)} + b, &\text{if} \, \sum(M) > 0\\
		0, &\text{otherwise ,}\\
		\end{cases}
	\end{align}
	
	where $W$ denotes the convolutional filter weights, $b$ its corresponding bias, $X$ the feature values for the current convolution, $M$ the corresponding mask and $1$ an array with the same size as $M$ filled with ones \cite{liu_image_2018}. This mask is then updated for each element according to \cite{liu_image_2018} as
	
	\begin{align}
		m' = \begin{cases}
		1, & \text{if} \sum(M) > 0 \\
		0, & \text{otherwise .} \\
		\end{cases}
	\end{align}
	
	So, with enough partial convolutions, the mask is eventually updated to incorporate only ones and the output image does not possess any empty pixels anymore. Decoding layers then additionally use the \textit{PyTorch} interpolation function to upsample the image successively back to its original scale.\\
	This whole process is then passed to the forward function of the neural network and conducted for all encoding and decoding layers using ReLU activation functions. Afterwards, the loss is calculated by separating the overall loss function into five parts, each for a different feature of the network output. On the one hand, the loss functions for the infilled holes and originally unmasked data points are weighted individually. On the other hand, their features are also added separately to form the total loss of the output image. Furthermore, a value for the total variation loss is added to account for the variation between neighboring image cells. The individual components are always calculated as mean absolute error (MAE) between the network output and the ground truth of the training data. The weights of the individual loss components are determined manually by tryouts without any specific hyperparameter training. Finally, the network weights are updated using \textit{PyTorch's} backpropagation and optimization functions according to the overall loss function. 
	
	\subsection{Neural Network Setup and Parameters}
	
	The U-net conceptually described in section \ref{method:unet} can be setup in a number of different ways depending on the exact purpose of the training and the input data. In general, the network is trained to reconstruct each slice of the temperature profile separately and then reassembled after the reconstruction. As the aim is to calculate the upper 700 m ocean heat content, the first 20 depth steps of the assimilation training data are selected, which represent temperature fields up to a depth of 700 m. 20 different networks then train to reconstruct the individual depth layers. \\
	The hyperparameters of the training, such as network depth or loss parameters can be adjusted previous to each training instance. The described setup with four encoding and four pooling layers is taken as a default setting, because until now, a further deepening of the neural network slows down calculations without providing significantly better results. Furthermore, the time period as well as the spatial extension in the form of longitudes and latitudes of the input data can be individually determined for each training process. In regard to the time period of the training data, two different arguments can be adjusted. On the one hand, the origin of the data itself can be adjusted, with default options from he Argo era, pre-Argo era or the entire assimilations time period. On the other hand, the time of origin for the observational patterns, with which the assimilation data is masked during the training process can be changed. Thus, it is possible for the network to learn Argo-era assimilation patterns while also learning to adjust to pre-Argo era observational density. \\
	The remaining parameters, such as learning rate, batch size and image size, are adopted from the original network set up by \cite{kadow_artificial_2020} and can be viewed with the remainder of the code under \url{https://github.com/slenta/Master-Arbeit}.
	
	\subsection{Data Preprocessing}
	
	The U-net described in section \ref{method:unet} takes three input images in the form of a ground truth assimilation, a binary mask and a masked out assimilation image formatted as an hdf5 file. In order feed the NetCDF4 files containing the assimilation data into the network, they have to be preprocessed to fit the given format. At first, the assimilation's temperature profiles are converted to anomalies using the assimilation's average monthly climatology. For this calculation only the time period from 2004 -- 2020 is considered, as this work assumes that the assimilation possesses significantly less uncertainty during the well observed Argo era. Then, EN4 observations from a selected input time are converted to binary masks and multiplied with the data assimilation model data. All three, the masks, the masked data as well as the data assimilations are then stored in an hdf5 file. From there, they can be loaded, split up into batches and fed into the network using a class with the \textit{PyTorch} dataset module.\\
	After the training process, when reconstructing a dataset with monthly observations for a certain timespan, the U-net returns outputs with fully reconstructed monthly temperature profiles.
	
	
	\section{Reconstructing North Atlantic Ocean Heat Content}
	
	The aim of this work is to develop an improved uncertainty estimate on North Atlantic ocean heat content, specifically in the time before the availability of Argo data. To do this, I propose to divide the available 16 member ensemble of monthly data assimilation runs into two distinct eras. The first part will be called the pre-Argo era, which starts with the beginning of the assimilation runs in 1958 and ends once the Argo project has reached close to global coverage of subsurface temperature observations in 2004. The timespan afterwards, once the Argo floats have reached near global observation coverage, will be called the Argo era and ends with the end of the assimilation runs in October 2020. \\
	
	To assess, if convolutional U-nets can be considered a valid alternative to EnKF data assimilation, the neural network will be trained on data assimilations during the well observed Argo era (2004 -- 2020) in order to learn on assimilations, which posses lower uncertainty and have integrated more observational data. After the network has been fully trained, visible at the progression of its training curve, it can be given masked temperature profiles for reconstruction. At first, the network will reconstruct the Argo era masked assimilations from one of the held back ensemble members. Each month of assimilation data will be masked with the corresponding observational pattern from that time. Given the high observational density during the Argo era, this reconstruction can be considered the first step for the neural network towards the reproduction of the assimilation OHC estimate. This approach can then be extended to the pre-Argo era, if the neural network is able to reproduce the assimilation ensemble within its uncertainty range during the Argo era. Analyzing the reconstructed OHC timeseries as well as maps from individual timesteps, the machine learning reconstructions can also be used to discuss inconsistencies and uncertainties within the assimilation. Differences between both OHC estimates can then be taken as indicator for uncertainties that are not mirrored within the assimilation's ensemble spread. Finally, I will examine, whether the neural network is able to directly reconstruct observations in the pre-Argo and Argo era as well as more recent observations it has not trained on within the assimilations uncertainty range. The results will then be evaluated in regard to the neural network's potential to reproduce or replace EnKF assimilations. Additionally, its implications on a new uncertainty estimate on the EnKF assimilations OHC outside of its ensemble spread will be considered. \\

	
	
\end{document}